\documentclass[11pt,twoside]{scrbook}

\usepackage[utf8]{inputenc}
\usepackage[T1]{fontenc}
\usepackage{lmodern}
\usepackage{booktabs}
\usepackage{natbib}

\title{Configuring a Book from First Principles}
\author{Noor Fatima}
\date{\today}

\begin{document}

\frontmatter
\maketitle

\addchap{Preface}
This book uses \texttt{scrbook}, the KOMA-Script book class, which favours
sensible defaults derived from classical page-design ratios over the fixed
choices of the standard \texttt{book} class.

\tableofcontents

\mainmatter

\chapter{Page Design}

\section{Why the Margins Look Different}

KOMA-Script computes its page geometry from the type area and paper size
rather than hardcoding margin lengths \citep{knuth1984texbook}. This produces
a different default layout from \texttt{book}, though every dimension remains
adjustable via the \texttt{typearea} package that \texttt{scrbook} loads.

\section{Headings and Chapter Marks}

Table~\ref{tab:levels} lists the sectioning commands \texttt{scrbook}
inherits, unchanged in name from the standard classes.

\begin{table}[h]
  \centering
  \begin{tabular}{ll}
    \toprule
    Command & Level \\
    \midrule
    \verb|\part|    & -1 \\
    \verb|\chapter| & 0 \\
    \verb|\section| & 1 \\
    \bottomrule
  \end{tabular}
  \caption{Sectioning commands available in \texttt{scrbook}.}
  \label{tab:levels}
\end{table}

\chapter{Fonts and Bibliography}

\section{Font Selection}

\texttt{scrbook} leaves font selection to the document, so \texttt{lmodern}
is loaded here explicitly rather than assumed.

\section{Closing the Bibliography}

A standard \texttt{natbib} bibliography closes out the back matter, exactly
as it would under \texttt{book} \citep{knuth1984texbook}.

\bibliographystyle{plainnat}
\bibliography{references}

\end{document}
